\def\level{Première}  % Classe à droite
\def\course{NSI}
\def\chaptername{Algèbre de Boole} % Titre au centre
\def\docname{RD02}        % Identifiant à gauche

\pagestyle{cours_FP}
\makeheader{} % Affiche le bandeau


\begin{minipage}{0.18\linewidth}
	\begin{center}
		\includegraphics[width = 0.95\linewidth]{George_Boole_color.jpg}

George Boole

(1815-1864)
	\end{center}
	
\end{minipage}\hfill{}
\begin{minipage}{0.8\linewidth}
 Les machines traitent et mémorisent l'information au moyen de circuits logiques binaires:

leurs entrées et sorties se caractérisent exclusivement par deux états :

\begin{itemize}
	\item l'état logique bas
	\item l'état logique haut

\end{itemize}

Ceci s'explique par la technologie employée: les microprocesseurs sont constitués d'une multitude de composants électroniques que l'on appelle des transistors et qui ne peuvent prendre que deux états, bloqué ou saturé, et se comportent comme des interrupteurs.

Ce sont des circuits logiques, le plus souvent à base de transistors, qui réalisent toutes les opérations dans les processeurs des machines.
\end{minipage}


\section{Variables booléennes}

\begin{definition}[Variable booléenne]{}
	\begin{itemize}
		\item     une variable booléenne est une variable qui ne peut prendre que deux valeurs : \textbf{vrai} ou \textbf{faux}.
		\item     les variables booléennes sont utilisées pour modéliser des situations à deux états, par exemple : allumé/éteint, ouvert/fermé, 1/0, \ldots{}
	\end{itemize}
\end{definition}

Au sein d'une machine, la valeur booléenne est donnée par l'état d'une grandeur physique. (intensité du courant, différence de tension etc.)

\section{Opérateurs booléens}

Les calculs sur les variables booléennes sont réalisés à l'aide d'opérateurs logiques. Les trois opérateurs logiques de base sont:

\subsubsection*{Opérateur NON / NOT (négation)}

\begin{minipage}{0.80\linewidth}
	L'opérateur NON s'applique sur un seul opérande et se note avec une barre au dessus de l'opérande. Il permet d'inverser la valeur d'une variable booléenne. Ainsi, si $A$ est une variable booléenne, alors $\overline{A}$ est la négation de $A$.
\end{minipage}\hfill{}
\begin{minipage}{0.18\linewidth}
	\begin{flushright}
		\renewcommand{\arraystretch}{1.2}
	\begin{tabular}{cc}
		\hline
		$A$ & $\overline{A}$ \\
		\hline
		False & True \\
		True & False \\
		\hline
		
	\end{tabular}
	\end{flushright}
	
\end{minipage}

\subsubsection*{Opérateur ET / AND (conjonction)}

\begin{minipage}{0.80\linewidth}
	L'opérateur ET s'applique sur deux opérandes et se note avec un point ou une croix entre les deux opérandes. Il permet de vérifier que les deux variables booléennes sont vraies en même temps. Ainsi, si $A$ et $B$ sont des variables booléennes, alors $A \wedge B$ est la conjonction de $A$ et $B$.
\end{minipage}\hfill{}
\begin{minipage}{0.18\linewidth}
	\begin{flushright}
		\renewcommand{\arraystretch}{1.2}
	\begin{tabular}{ccc}
		\hline
		$A$ & $B$ & $A \cdot B$ \\
		\hline
		False & False & False \\
		False & True & False \\
		True & False & False \\
		True & True & True \\
		\hline
		
	\end{tabular}
	\end{flushright}
	
\end{minipage}

\subsubsection*{Opérateur OU / OR (disjonction)}

\begin{minipage}{0.80\linewidth}
	L'opérateur OU s'applique sur deux opérandes et se note avec un \verb|+| entre les deux opérandes. Il permet de vérifier que au moins une des deux variables booléennes est vraie. Ainsi, si $A$ et $B$ sont des variables booléennes, alors $A + B$ est la disjonction de $A$ et $B$.
\end{minipage}\hfill{}
\begin{minipage}{0.18\linewidth}
	\begin{flushright}
		\renewcommand{\arraystretch}{1.2}
	\begin{tabular}{ccc}
		\hline
		$A$ & $B$ & $A + B$ \\
		\hline
		False & False & False \\
		False & True & True \\
		True & False & True \\
		True & True & True \\
		\hline
		
	\end{tabular}
	\end{flushright}
	
\end{minipage}

\pagebreak

\begin{propriete}[Notations]{}
 Par abus de langage, on notera:
\begin{itemize}
	\item \verb|1| pour \verb|True|
	\item \verb|0| pour \verb|False|
\end{itemize}

On peut aussi utiliser les notations suivantes pour les opérateurs logiques:
\begin{itemize}
	\item $A \cdot B$ ou $AB$ pour $A \wedge B$
	\item $A + B$ pour $A \vee B$
	\item $\overline{A}$ pour $\neg A$
\end{itemize}
\end{propriete}

\begin{exercice_cours}[Montrer que l'opérateur OU peut s'écrire avec des opérateurs NON et ET]{}

La méthode est de prouver que $A + B$ et $\overline{\overline{A} \cdot \overline{B}}$ ont la même table de vérité.

\begin{multicols}{2}
\renewcommand{\arraystretch}{1.3}
\begin{tabular}{ccc}
	\hline
	$A$ & $B$ & $A + B$ \\
	\hline
	False & False & False \\
	False & True & True \\
	True & False & True \\
	True & True & True \\
	\hline
	

\end{tabular}

\columnbreak
\renewcommand{\arraystretch}{1.3}
\begin{tabular}{cc|p{1.2cm}|p{1.2cm}|p{1.2cm}|p{1.2cm}|}
	\hline
	$A$ & $B$ & $\overline{A}$ & $\overline{B}$ & $\overline{A} \cdot \overline{B} $ & $\overline{\overline{A} \cdot \overline{B}}$ \\
	\hline
	False & False &  &  &  &  \\
	False & True &  &  &  &  \\
	True & False &  &  &  &  \\
	True & True &  &  &  &  \\
	\hline
\end{tabular}
\end{multicols}

\end{exercice_cours}

\begin{exercice_cours}[Dresser la table de vérité de $A \cdot (\overline{B} + C)$]{}

\renewcommand{\arraystretch}{1.2}
\begin{tabular}{ccc|p{1.5cm}|p{1.5cm}|p{1.5cm}|p{1.5cm}|p{1.5cm}|p{1.5cm}|p{1.5cm}|p{1.5cm}|}
	\hline
	$A$ & $B$ & $C$ & $\overline{B}$ & $\overline{B} + C$ & $A \cdot (\overline{B} + C)$ \\
	\hline
	False & False & False &  &  &  \\
	False & False & True &  &  &  \\
	False & True & False &  &  &  \\
	False & True & True &  &  &  \\
	True & False & False &  &  &  \\
	True & False & True &  &  &  \\
	True & True & False &  &  &  \\
	True & True & True &  &  &  \\
	\hline
\end{tabular}

\end{exercice_cours}

\subsection*{Opérateur OU exclusif / XOR (disjonction exclusive)}

\begin{minipage}{0.80\linewidth}
	L'opérateur OU exclusif s'applique sur deux opérandes et se note avec un $\oplus$ entre les deux opérandes. Il permet de vérifier que les deux variables booléennes sont différentes. Ainsi, si $A$ et $B$ sont des variables booléennes, alors $A \oplus B$ est la disjonction exclusive de $A$ et $B$.
\end{minipage}\hfill{}
\begin{minipage}{0.18\linewidth}
	\begin{flushright}
		\renewcommand{\arraystretch}{1.2}
	\begin{tabular}{ccc}
		\hline
		$A$ & $B$ & $A \oplus B$ \\
		\hline
		False & False & False \\
		False & True & True \\
		True & False & True \\
		True & True & False \\
		\hline
		
	\end{tabular}
	\end{flushright}
	
\end{minipage}

\begin{exercice_cours}[Montrer que l'opérateur OU exclusif peut s'écrire $A \oplus B = A \cdot \overline{B} + \overline{A} \cdot B$]{}

\begin{tabular}{cc|p{1.5cm}|p{1.5cm}|p{1.5cm}|p{1.5cm}|p{1.5cm}|p{1.5cm}|p{1.5cm}|p{1.5cm}|}
	\hline
	$A$ & $B$ & $A \oplus B$ & $A \cdot \overline{B}$ & $\overline{A} \cdot B$ & $A \cdot \overline{B} + \overline{A} \cdot B$ \\
	\hline
	False & False &  &  &  &  \\
	False & True &  &  &  &  \\
	True & False &  &  &  &  \\
	True & True &  &	  &  &  \\
	\hline
\end{tabular}

\end{exercice_cours}

\pagebreak

\section{Propriété des opérateurs}

\begin{propriete}[Element neutre]{}
	\begin{itemize}
		\item L'élément neutre de l'addition est 0: $A + 0 = A$
		\item L'élément neutre de la multiplication est 1: $A \cdot 1 = A$
	\end{itemize}
\end{propriete}

\begin{propriete}[Element absorbant]{}
	\begin{itemize}
		\item L'élément absorbant de l'addition est 1: $A + 1 = 1$
		\item L'élément absorbant de la multiplication est 0: $A \cdot 0 = 0$
	\end{itemize}
\end{propriete}

\begin{propriete}[Elements complémentaires et incompatibles]{}
	\begin{itemize}
		\item $A + \overline{A} = 1$
		\item $A \cdot \overline{A} = 0$
	\end{itemize}
\end{propriete}

\begin{propriete}[Loi de Morgan]{}
	\begin{itemize}
		\item $\overline{A + B} = \overline{A} \cdot \overline{B}$
		\item $\overline{A \cdot B} = \overline{A} + \overline{B}$
	\end{itemize}
\end{propriete}

\section{Algèbre de Boole et informatique}

\begin{minipage}[t]{0.49\linewidth}
	\subsubsection*{Les opérateurs booléens (ou logiques)}

\begin{tabular}{|c|c|c|c|}
	
	\textbf{Opération} & \textbf{Logique} & \textbf{Python} & \textbf{C/C++/JS/PHP} \\
	\hline
	FAUX & 0 & False & false \\
	VRAI & 1 & True & true \\
	NON & $\overline{A}$ & not A & !A \\
	ET & $A \cdot B$ & A and B & A \verb|&&| B \\
	OU & $A + B$ & A or B & A \verb+||+ B \\
	\hline
\end{tabular}


\end{minipage}\hfill
\begin{minipage}[t]{0.49\linewidth}
	\subsubsection*{Opérateurs booléens binaires}

\begin{tabular}{|c|c|c|}
	
	\textbf{Opération} & \textbf{Logique} & \textbf{Python/C/C++/JS/PHP} \\
	\hline
	ET bit à bit & $A \cdot B$ & A \verb|&| B  \\
	OU bit à bit & $A + B$ & A \verb+|+ B  \\
	XOR bit à bit & $A \oplus B$ & A \verb|^| B  \\
	NON bit à bit & $\overline{A}$ & \verb|~|A  \\
	\hline
\end{tabular}

\end{minipage}

\begin{exercice_cours}[Quelle est la valeur des expression suivantes en Python?]{}

\begin{itemize}
	\item \verb|True and False or True| \dotfill
	\item \verb|True or False and True| \dotfill
	\item \verb|not True| \dotfill
	\item \verb|True and not False| \dotfill
	\item \verb|False or not False| \dotfill
\end{itemize}
\end{exercice_cours}